% profiles/indiana-university-bloomington.tex — Indiana University-Bloomington
% ============================================================================
% source: https://graduate.indiana.edu/academic-requirements/thesis-dissertation/doctoral-guide/formatting.html
%         (The Graduate School Bloomington, Doctoral Dissertation Guide,
%          "Formatting")
% source: https://graduate.indiana.edu/academic-requirements/thesis-dissertation/doctoral-guide/sections.html
%         (the same guide, "Required & Optional Sections")
% source: https://graduate.indiana.edu/academic-requirements/thesis-dissertation/masters-guide/sections.html
%         (Master's Thesis Guide, "Required & Optional Sections")
% source: https://graduate.indiana.edu/doc/phdtitlepage.docx
% source: https://graduate.indiana.edu/doc/phdacceptancepage.docx
% source: https://graduate.indiana.edu/doc/masters-title-page.docx
% source: https://graduate.indiana.edu/doc/masters-acceptance-page.docx
%         (the four specimens the guides link as "an example of the ... page")
% accessed: 2026-09-09
% checked:  texlib-thesis-profile-round
%
% VERIFIED and set here: geometry, spacing, the title page, and the acceptance
% page. Both degree levels are covered; the guides differ only in the degree
% named and in the committee heading.
% NOT set:
%   * \thesissetchapteropening — none prescribed.
%   * every checkable accessibility declaration — see the accessibility note.
%
% CURRENCY. The pages carry no revision date. They are plainly maintained (they
% reference ProQuest submission and the current Graduate Recorder address), and
% no superseded edition was found. Treat the substance as current but note that
% nothing here is dated.
%
% THE ACCEPTANCE PAGE IS REQUIRED AND IS FILED UNSIGNED. Indiana is unusual in
% wanting the signature RULES in the document with nothing on them:
%
%     "Your Research Committee members' names and signature lines follow and are
%      fully right-justified."
%     "The page remains unsigned for inclusion in your dissertation."
%
% So the rules below are correct as drawn, and a filer must not replace this
% page with a scan of a signed one. That is the opposite of Mount Sinai and
% Arizona, and matches UC Davis's arrangement.
%
% THE TYPEFACE LIST IS A RECOMMENDATION, so the class passes: the six faces
% Indiana names (Arial, Bookman Old Style, Calibri, Cambria, Lucida Bright,
% Times New Roman) are "recommended for the ease with which they convert to a
% PDF", not required. Size is: "Font size should be either 11 or 12 point for
% the entire document", with the title page allowed up to 16 point and footnotes
% and table content down to 10. The class's 12pt is in range.
%
% ONE PROHIBITION WORTH KNOWING, and one this class was never going to break:
% "The IU Seal or Branding should not be used on any portion of the
% dissertation. These items may be used only with the written permission of the
% university." No insignia is drawn here.
%
% PAGE NUMBERING, which no profile field can express: "The title page counts as
% page i but does not bear a number. Begin the actual numbering with the
% acceptance page as page ii, and continue with lowercase Roman numerals until
% the start of the actual body", then Arabic 1 onward. The vita at the end
% carries no page number at all, and "Running heads are not used in dissertation
% submissions. Please limit the content of your header and footer space to the
% page number only." A document filed here needs \frontmatter before the title
% page and its own handling of the unnumbered title and vita pages.

\thesisinstitution{Indiana University}

% "Top, Right, and Bottom margins must be one inch. If the dissertation will
% only be electronically accessed, a one-inch left margin is acceptable. The
% left margin should be one inch if the dissertation will be bound in paper form
% by ProQuest. If using a bindery other than ProQuest, please consult with the
% bindery about the size of the left margin needed for their binding process."
%
% One inch on all four sides is therefore right for both the electronic filing
% and a ProQuest-bound copy. The only case Indiana leaves open is a third-party
% bindery, which it sends the student to ask -- so a wider left margin is a
% decision for that student and their bindery, not a figure this profile can
% pick. "These margin requirements apply to all materials included in the
% dissertation, including figures, tables, maps, plates, etc."
\thesissetgeometry{left=1in, right=1in, top=1in, bottom=1in}

% "The material should be double-spaced. Long quotations within the text should
% be typed single-spaced with wider margins." Not a choice here.
\thesissetspacing{double}

% --- Accessibility -----------------------------------------------------------
% Both guides -- Formatting and Required & Optional Sections, at both degree
% levels -- were read end to end, along with the Thesis & Dissertation index.
% None of them contains the words accessible, accessibility, alt text, tagged,
% WCAG, PDF/UA, PDF/A, screen reader or Title II, and none states a
% document-language or heading-structure rule.
%
% IU Libraries publishes guidance on accessible PDFs. That is a library research
% guide, not the Graduate School's filing requirements, so under RESEARCHING.md
% it is not a source for this file and nothing from it is encoded here.
\thesisaccessibilityrequirement{Not published}
	{The Graduate School Bloomington's doctoral and master's guides state no
	 accessibility requirement. All of Formatting, Required \& Optional Sections
	 and the Thesis \& Dissertation index were read in full on the access date.
	 This records that the Graduate School asks for nothing on accessibility,
	 not that nobody looked.}

% --- Profile-local metadata ----------------------------------------------------
% Indiana's two front pages carry two DIFFERENT dates, and conflating them is
% the easy mistake: \thesis@date is the title page's, and the guide is explicit
% that it is not the defence -- "(The month and year is the date when all
% requirements have been satisfied -- this is not necessarily the month in which
% you defend.)" The acceptance page carries the defence date, in full, and that
% is what \iudefensedate supplies.
\newcommand{\iu@defense}{}
\newcommand{\iudefensedate}[1]{\renewcommand{\iu@defense}{#1}}
% The department or school named on the title page, if it is not what
% \gradprogram already says.
\newcommand{\iu@unit}{}
\newcommand{\iuunit}[1]{\renewcommand{\iu@unit}{#1}}
\newcommand{\iu@theunit}{\ifx\iu@unit\@empty\thesis@program\else\iu@unit\fi}
% "Doctoral Committee" or "Master's Committee", chosen from \doctype.
\newcommand{\iu@ifdoctoral}[2]{%
	\def\iu@tmp{dissertation}%
	\ifx\thesis@docnoun\iu@tmp #1\else #2\fi
}

% --- Title page ----------------------------------------------------------------
% "The title page should be a separate page and no longer than one page. All
% content on this page should be centered horizontally and vertically. Keep in
% mind that the title must be able to fit on the spine of a bound manuscript."
%
% The doctoral specimen, whose first line is an instruction to the reader rather
% than page content ("TITLE CENTERED, ALL IN CAPITAL LETTERS:"):
%
%     THE IMPORTANCE OF KEY WORDS IN THE SUCCESSFUL INDEX
%     Author's Name
%     Submitted to the faculty of the Graduate School in partial fulfillment of
%     the requirements
%     for the degree Doctor of Philosophy
%     in the Department (or School) of ______________ ,
%     Indiana University
%     Month Year
%
% The master's specimen is the same page with "for the degree Master of (Arts or
% Sciences)" and the department and university on one line, so \thesis@degree
% carries the difference and no branch is needed. Note "for the degree", with no
% "of" -- the acceptance page says "for the degree of", and they really do
% differ.
%
% The title page bears no page number, though it counts as page i.
\renewcommand{\thesistitlepage}{%
	\loadgeometry{nopagenumber}%
	\begin{titlepage}%
		\thispagestyle{empty}\centering\singlespacing
		\vspace*{\fill}
		\thesis@tagtitle{Title}{\MakeUppercase{\@title}}\par
		\vspace{4\baselineskip}
		\@author\par
		\vspace{4\baselineskip}
		Submitted to the faculty of the Graduate School in partial fulfillment
		of the requirements\\
		for the degree \thesis@degree\\
		in the \iu@theunit,\\
		Indiana University\\
		\thesis@date\par
		\vspace*{\fill}
	\end{titlepage}%
	\loadgeometry{normal}%
}

% --- Acceptance page -----------------------------------------------------------
% "This page confirms the committee's approval and acceptance of your
% dissertation. The acceptance page should be a separate page and no longer than
% one page." The layout is prescribed in four sentences, each honoured here:
%
%   - "The Acceptance statement should be centered at the top of the page:
%     Accepted by the Graduate Faculty, Indiana University, in partial
%     fulfillment of the requirements for the degree of Doctor of Philosophy."
%   - "The words 'Doctoral Committee' immediately follow the Acceptance
%     statement and are left-justified." (Master's: "Master's Committee".)
%   - "Your Research Committee members' names and signature lines follow and are
%     fully right-justified. Behind each name, place the appropriate
%     post-nominal initials for that individual (Ph.D., M.F.A., O.D., etc.)"
%   - "After the names, should be the date of your defense, left-justified."
%
% ONE DISAGREEMENT BETWEEN PROSE AND SPECIMEN, left for a reviewer rather than
% split: the master's specimen puts the defence date immediately after the
% "Master's Committee" heading and BEFORE the signature lines, while the prose
% at both degree levels says it comes after the names. The doctoral specimen
% agrees with the prose. This page follows the prose, per RESEARCHING.md.
%
% Post-nominal initials belong to the NAME, not to the class's role argument:
% pass \committeemember{Robin Reader, Ph.D.}{} and leave the second argument
% empty. \thesis@memberline is replaced because the class's default rule is
% centred and 3in wide, where Indiana wants the rule and the name both flush
% right.
\renewcommand{\thesisapprovalpage}{%
	\clearpage\thispagestyle{empty}%
	\loadgeometry{nopagenumber}\singlespacing
	\thesis@tagtitle{H1}{{\color{white}\tiny Acceptance Page}}\par
	\begin{center}
		Accepted by the Graduate Faculty, Indiana University, in partial
		fulfillment of the requirements for the degree of \thesis@degree.
	\end{center}
	\vspace{2\baselineskip}
	\noindent\iu@ifdoctoral{Doctoral Committee}{Master's Committee}\par
	\vspace{\baselineskip}
	\begingroup
	\renewcommand{\thesis@memberline}[2]{%
		\vspace{2\baselineskip}
		{\raggedleft\rule{3.4in}{0.4pt}\par
		 ##1\ifx\relax##2\relax\else, ##2\fi\par}%
	}%
	\thesis@committee
	\endgroup
	\vspace{3\baselineskip}
	\noindent\ifx\iu@defense\@empty\thesis@date\else\iu@defense\fi\par
	\vspace*{\fill}
	\loadgeometry{normal}\doublespacing\clearpage
}
